%\noindent 
在前几部分中，我们主要从技术视角考察智能体大语言模型：架构、训练流程、推理策略、工具使用以及在各个领域的应用。在本部分中，我们换个角度。我们不再追问智能体大语言模型如何工作、如何构建，而是追问它们在{\em 社会性}上如何表现。%Previous parts studied how agents act, this part studies how they act {\em together}. 
随着智能体大语言模型越来越多地被部署到社会情境中（作为助手、协作者、导师等），它们的行为倾向开始具有实际意义。例如，一个智能体在策略情境中倾向于合作还是背叛、它是否会对礼貌用语作出反应、它能否在多智能体模拟中追踪另一个智能体的信念状态——这些都有现实后果。  

本部分讨论的大多数研究，都建立在最初为研究人类与其他动物而发展出的方法之上，涉及策略行为、社交礼仪、共情、人格等特质。每章末尾，我们都会就一个普遍性问题做简短反思：这类方法能否有效地应用于 AI。  

在下一部分、也是最后一部分中，我们将更深入地探讨大语言模型的内部运作，以及各种形式的认知、乃至意识可能涌现的问题。
% The preceding four parts showed us how to create LLM Agents that reason, act, and interact with the world. This poses the question how such agents will behave in different contexts. Will they cooperate or compete, lead or follow? Is their behavior stable and predictable, or rather dynamic in novel ways? And what kind of behavior will emerge in larger systems of interacting LLM agents? The first chapter of this part will study different kinds of role-based and experimentally controlled behavior. The second chapter discusses emergent behavior in more open-ended settings. The final content chapter goes deeper, and will dive into parallels and differences between behaviour and cognition of LLM agents and humans. Finally, we will discuss a practical simulation of a society of Agents.
